\scanpage{98}

\[
\begin{aligned}
 &\text{L'application }\quad\lambda\longmapsto\lambda.i_0\qquad Z\longrightarrow\Pi\\
 &\text{est bijective ; plus précisément,}\\
 &\text{si on considère de \unsure{$\simeq$|$=$} l'application}\\
 &\lambda\longmapsto\lambda.i,\ \text{pour tt choix de }i\in\Pi,\\
 &\text{on trouve que l'ens des structures}\\
 &\text{2-polyédrales }\Pi\text{ sur }S\text{ est un torseur}\\
 &\text{sous }Z.
\end{aligned}\tag*{(59)}
\]

Il faudrait maintenant expliciter l'iso

\[
Z\simeq Z'_\Lambda\simeq\Lambda^*/\{\pm 1_\Lambda\}\tag*{(60)}
\]

en termes de la seule structure combinatoire d'un des $P_i$, soit $P_{i_0}$. Regardons pour ceci, pour $\lambda\in Z$ fixé, et $s\in S$ fixé, l'application

\[
\left\{\begin{aligned}
 \operatorname{Et}_{P_{i_0}}(s)&\overset{\sim}{\longrightarrow}\operatorname{Et}_{P_{\underbrace{\lambda.i_0}_{i_1}}}(s)\\
 t&\longmapsto\lambda.t
\end{aligned}\right.\tag*{(61)}
\]

Notons que les deux mbres sont des torseurs sous les groupes

\[
\left\{\begin{aligned}
 &\deleted{\unsure{\mathbb Z/n\mathbb Z}}\\
 &\Gamma_{i_0}=\Lambda\wedge_{\mu_2}\mu(i_0)\qquad\text{resp }\Gamma_{i_1}=\Lambda\wedge_{\mu_2}\mu(i_1)\\
 &\text{où }\mu(i)=\text{ens des orientations de }P_i
\end{aligned}\right.\tag*{(62)}
\]
% archive sheet number: 97
